% META: {"id": "TSPE-DERCONV-CR-013", "chapitre": "TSPE-DERIVATION-CONVEXITE", "type_objet": "cours", "capacites_codes": ["C4"], "sources_inspiration": [], "status": "approved"}

\section{Esquisse de courbe a partir de f, f' et f''}

% ============================================================
% STRATE 1 — Liens entre f, f', f''
% ============================================================

\propriete{
Soit $f$ une fonction deux fois dérivable sur un intervalle $I$.

\begin{itemize}
  \item $f'(x) > 0$ : $f$ est strictement croissante.
  \item $f'(x) < 0$ : $f$ est strictement décroissante.
  \item $f''(x) > 0$ : $f'$ est strictement croissante, la courbe de $f$ « tourne vers le haut » (convexe).
  \item $f''(x) < 0$ : $f'$ est strictement décroissante, la courbe de $f$ « tourne vers le bas » (concave).
\end{itemize}
}

\definition[1]{
Un \textbf{point d'inflexion} de la courbe de $f$ est un point ou la courbe change de convexité (passage de convexe a concave ou inversement). En ce point, la tangente traverse la courbe.
}

\propriete{
Si $f$ est deux fois dérivable et si $f''$ s'annule en $a$ en changeant de signe, alors le point $(a, f(a))$ est un point d'inflexion de la courbe.
}

\subsection{Construire une esquisse}

\propriete{
Pour esquisser la courbe de $f$ a partir des tableaux de $f'$ et $f''$ :
\begin{enumerate}
  \item Placer les extremums (la ou $f'$ s'annule en changeant de signe) et les points d'inflexion (la ou $f''$ s'annule en changeant de signe).
  \item Dans chaque intervalle, tracer la courbe en respectant :
    \begin{itemize}
      \item le sens de variation (donne par le signe de $f'$),
      \item la courbure (donnee par le signe de $f''$).
    \end{itemize}
  \item Tracer les tangentes aux points d'inflexion (la courbe les traverse).
\end{enumerate}
}

\exemple{
Soit $f(x) = x^3 - 3x$.

$f'(x) = 3x^2 - 3 = 3(x-1)(x+1)$.

$f''(x) = 6x$.

$f'$ s'annule en $x = -1$ (maximum local) et $x = 1$ (minimum local).

$f''$ s'annule en $x = 0$ et change de signe : point d'inflexion en $(0, 0)$.

Pour $x < 0$ : $f'' < 0$ (concave). Pour $x > 0$ : $f'' > 0$ (convexe).
}

% ============================================================
% STRATE 2 — Erreurs frequentes
% ============================================================

\erreurFrequente{
\textbf{Confondre extremum et point d'inflexion.} Un extremum est un point ou $f'$ s'annule en changeant de signe. Un point d'inflexion est un point ou $f''$ s'annule en changeant de signe. Ce ne sont pas les mêmes objets.
}

\erreurFrequente{
\textbf{Croire que $f''(a) = 0$ suffit pour un point d'inflexion.} Il faut que $f''$ change de signe en $a$. Par exemple, $f(x) = x^4$ vérifie $f''(0) = 0$ mais $(0, 0)$ n'est pas un point d'inflexion (la fonction reste convexe).
}

\margeAppui{
\textbf{Resume visuel :}

$f'' > 0$ : courbe en forme de bol (convexe, « souriante »).

$f'' < 0$ : courbe en forme de chapeau (concave, « triste »).

Point d'inflexion : passage de l'un a l'autre.
}

% BEGIN-VERIFY
% from sympy import symbols, diff, solve, Rational
%
% x = symbols('x', real=True)
%
% # Un point d'inflexion exige que f'' s'annule EN CHANGEANT DE SIGNE.
% f = x**3
% assert solve(diff(f, x, 2), x) == [0]
% assert diff(f, x, 2).subs(x, -1) < 0 < diff(f, x, 2).subs(x, 1)
% # La tangente en 0 est l'axe des abscisses, et la courbe la traverse.
% assert f.subs(x, -1) < 0 < f.subs(x, 1)
%
% # Le contre-exemple decisif : f'' peut s'annuler SANS changement de signe,
% # et il n'y a alors pas de point d'inflexion. Sur x^4 :
% g = x**4
% assert solve(diff(g, x, 2), x) == [0]
% assert diff(g, x, 2).subs(x, -1) > 0 and diff(g, x, 2).subs(x, 1) > 0
% # La courbe reste au-dessus de sa tangente en 0 des deux cotes.
% assert g.subs(x, -1) > 0 and g.subs(x, 1) > 0
%
% # Signe de f' et sens de variation, sur x^3 - 3x.
% h = x**3 - 3*x
% hp = diff(h, x)
% assert sorted(solve(hp, x)) == [-1, 1]
% assert hp.subs(x, -2) > 0 and hp.subs(x, 0) < 0 and hp.subs(x, 2) > 0
% assert h.subs(x, -2) < h.subs(x, -1)          # croissante avant -1
% assert h.subs(x, -1) > h.subs(x, 0)           # decroissante entre -1 et 1
% assert h.subs(x, 1) < h.subs(x, 2)            # croissante apres 1
% END-VERIFY
